\chapter{Démonstrations : le modèle de cascade}
\label{ann:cascade}

% HARNAIS-HORS-SCRIPT: chapitre de demonstrations de la cascade : memes constantes d'enonce ; les valeurs numeriques correspondantes sont verifiees au chapitre cascade, qui cite ses scripts.

Cette annexe démontre les propriétés énoncées à la section~\ref{sec:formalisation}. Les
notations sont celles du chapitre~\ref{chap:cascade} : $P=\{1,\dots,5\}$, $W\in[0,1]^{P\times P}$
avec $W_{jj}=0$, et la cascade indépendante de la définition~\ref{def:cascade}.

\section*{Forme réduite et descendance espérée (proposition~\ref{prop:reduite})}

\begin{proof}
Soit $\mathcal M_n(a,k)$ le nombre de marches dirigées de longueur $n$ de $a$ vers $k$ dans le
graphe pondéré. Une marche $a=v_0\to v_1\to\cdots\to v_n=k$ est vivante avec probabilité
$\prod_{t} W_{v_{t-1}v_t}$, ces arêtes étant indépendantes ; en sommant sur les marches,
$\mathbb E\,\mathcal M_n(a,k)=\sum \prod_t W_{v_{t-1}v_t}=(W^n)_{ak}$ par définition du produit
matriciel. Le nombre de fois où $k$ est atteint depuis $a$ en comptant \emph{chaque} chemin
vaut $\sum_{n\ge0}\mathcal M_n(a,k)$, d'espérance $\sum_{n\ge0}(W^n)_{ak}$. Cette série
matricielle est la série de Neumann de $W$ ; elle converge vers $(I-W)^{-1}$ si et seulement si
$\rho(W)<1$, auquel cas $\sum_{n\ge0}W^n=(I-W)^{-1}$.

L'indicatrice $\mathbf 1\{k\in S_a\}$ vaut $1$ dès qu'\emph{au moins un} chemin vivant relie
$a$ à $k$, donc $\mathbf 1\{k\in S_a\}\le \sum_{n\ge0}\mathcal M_n(a,k)$ presque sûrement. En
prenant l'espérance puis en sommant sur $k$,
\[
  \mathbb E\,|S_a|=\sum_k \mathbb P(k\in S_a)\le \sum_k\bigl((I-W)^{-1}\bigr)_{ak}
  =\bigl[(I-W)^{-1}\mathbf 1\bigr]_a .
\]
L'inégalité est stricte dès qu'un pilier est accessible par plusieurs chemins avec probabilité
positive, ce qui quantifie les collisions.
\end{proof}

\section*{Sous-criticité par normalisation (proposition~\ref{prop:leontief})}

\begin{proof}
Par construction~\eqref{eq:norm}, $W=(g/s_{\max})\,\mathrm{TRANS}\ge 0$ et, pour toute ligne
$j$, $\sum_k W_{jk}=(g/s_{\max})\sum_k \mathrm{TRANS}_{jk}\le (g/s_{\max})\,s_{\max}=g$, puisque
$s_{\max}$ est la plus grande somme de ligne de $\mathrm{TRANS}$. Pour une matrice à
coefficients positifs, le rayon spectral est majoré par toute norme d'opérateur subordonnée,
en particulier la norme infinie, qui est le maximum des sommes de lignes :
$\rho(W)\le\lVert W\rVert_\infty=\max_j\sum_k W_{jk}\le g<1$. La série de Neumann converge donc
et $(I-W)^{-1}=\sum_{n\ge0}W^n$ existe, à coefficients positifs.
\end{proof}

\section*{Loi exacte de l'ensemble atteint (proposition~\ref{prop:loi})}

\begin{proof}
Fixons $A\ni a$. L'événement $\{S_a=A\}$ est l'intersection de deux événements portant sur des
arêtes disjointes, donc indépendants : (i) aucune arête ne relie $A$ à son complémentaire (sans
quoi $S_a$ dépasserait $A$), et (ii) $a$ atteint tout $A$ par les seules arêtes internes à $A$.
L'événement (i) ne concerne que les arêtes $j\to k$ avec $j\in A$, $k\notin A$, chacune absente
avec probabilité $1-W_{jk}$ indépendamment ; sa probabilité est $B(A)=\prod_{j\in A,k\notin A}(1-W_{jk})$.
L'événement (ii) ne concerne que les arêtes internes à $A$ ; notons $R(A)$ sa probabilité.
Par indépendance, $\mathbb P(S_a=A)=R(A)\,B(A)$.

Pour $R(A)$, on partitionne selon le sous-ensemble $A'$ effectivement atteint par les arêtes
internes : $A'$ contient $a$, et $A'\subsetneq A$ est \emph{propre} si et seulement si $A$
n'est pas entièrement atteint. Or, conditionnellement à ce que l'ensemble atteint par les
arêtes internes soit exactement $A'$, aucune arête ne va de $A'$ vers $A\setminus A'$, événement
de probabilité $\prod_{j\in A',k\in A\setminus A'}(1-W_{jk})$ et indépendant de la réalisation
interne à $A'$. En sommant les probabilités de ces sous-ensembles propres et en passant au
complémentaire,
\[
  R(A)=1-\sum_{\substack{A'\subsetneq A\\ a\in A'}} R(A')\!\!\prod_{j\in A',k\in A\setminus A'}\!\!(1-W_{jk}),
\]
avec $R(\{a\})=1$. La récursion sur les sous-ensembles emboîtés est finie et détermine $R$ de
proche en proche.
\end{proof}

\section*{Monotonie stochastique (proposition~\ref{prop:mono})}

\begin{proof}
On construit un \emph{couplage}. Soit $(U_{jk})$ une famille de variables uniformes sur $[0,1]$
indépendantes, commune aux deux matrices. Pour une matrice $M$, déclarons l'arête $j\to k$
vivante si $U_{jk}<M_{jk}$ ; le graphe obtenu a la bonne loi (arête vivante avec probabilité
$M_{jk}$). Si $W\le W'$ terme à terme, alors $U_{jk}<W_{jk}\Rightarrow U_{jk}<W'_{jk}$ : toute
arête vivante sous $W$ l'est sous $W'$. Le graphe sous $W$ est donc un sous-graphe de celui sous
$W'$, et l'accessibilité étant croissante par ajout d'arêtes, $S_a(W)\subseteq S_a(W')$ presque
sûrement. Pour toute fonctionnelle $\varphi$ croissante pour l'inclusion,
$\varphi(S_a(W))\le\varphi(S_a(W'))$ p.s., d'où la domination stochastique et la monotonie de
l'espérance et des quantiles. Enfin, pour une perturbation antisymétrique $A$, l'entrée $W_{jk}$
croît quand $W_{kj}$ décroît : les deux effets jouent en sens contraire sur $S_a$, la monotonie
en $A_{jk}$ n'a plus lieu, et l'optimum sur le pavé $\{|A_{jk}|\le t\,S_{jk}\}$ doit être
cherché aux sommets sans garantie qu'ils le réalisent.
\end{proof}

\section*{Dépendance à l'ordre et non-identification (propositions~\ref{prop:ordre} et~\ref{prop:nonid})}

\begin{proof}
\emph{Ordre.} Soit $\sigma=(i_1,\dots,i_n)$ et $\sigma'$ une permutation de $\sigma$. Alors
$\ell(\sigma)=\ell(\sigma')$ pour toute paire si et seulement si le produit
$\prod_k \mathrm{TRANS}_{i_k i_{k+1}}$ ne dépend pas de l'ordre des facteurs disponibles, ce qui,
pour $n=2$, équivaut à $\mathrm{TRANS}_{ij}=\mathrm{TRANS}_{ji}$ pour tout $(i,j)$, soit la
symétrie de $\mathrm{TRANS}$. La contraposée donne le résultat : $\mathrm{TRANS}$ asymétrique
implique l'existence d'une paire dont l'ordre change la vraisemblance, donc la criticité n'est
pas une fonction de l'ensemble. Toute copule et tout score additif étant symétriques en leurs
arguments, aucun ne reproduit cette dépendance.

\emph{Non-identification.} La décomposition $W=\tfrac12(W+W^\top)+\tfrac12(W-W^\top)=S+A$ est
unique car $S$ est symétrique, $A$ antisymétrique, et cette somme directe est celle des
matrices en parties symétrique et antisymétrique. La co-occurrence observée d'une paire
$\{j,k\}$ est, par échange des rôles, une fonction de $W_{jk}+W_{kj}=2S_{jk}$ seulement ; de
même la statistique d'asymétrie du placebo compare $M$ à sa transposée et son espérance sous
permutation ne dépend que de $S$. Donc $W$ et $W^\top=S-A$ induisent les mêmes lois observables
de co-occurrence. Ils diffèrent par le signe de $A$, et la proposition~\ref{prop:mono} montre
qu'un tel changement modifie l'ensemble atteint, donc le SCR, en général : la direction est non
identifiée.
\end{proof}

\section*{Efficience de la valeur de Shapley}
\label{sec:ann-efficience-shapley}

Le chapitre~\ref{chap:resultats} attribue le surcoût de capital aux piliers par la valeur de
Shapley, et lit dans le tableau que la colonne des parts somme exactement à cent. Cette somme
n'est ni une coïncidence ni un artefact d'arrondi : c'est la propriété d'\textbf{efficience}
\citep{Shapley1953}, et elle se démontre en quelques lignes. Elle est reproduite ici parce que
c'est elle, et elle seule, qui autorise à présenter une attribution comme une \emph{partition} du
surcoût et non comme une pondération arbitraire.

\begin{proposition}[Efficience]
\label{prop:efficience}
Soit $v$ une fonction caractéristique sur $\{1,\dots,n\}$ avec $v(\emptyset)=0$, et
\[
  \varphi_i \;=\; \frac1n \sum_{u \subseteq -\{i\}} \binom{n-1}{|u|}^{-1}
                  \bigl[v(u \cup \{i\}) - v(u)\bigr].
\]
Alors $\sum_{i=1}^{n} \varphi_i = v(\{1,\dots,n\})$.
\end{proposition}

\begin{proof}
La démonstration repose entièrement sur la réécriture du poids. Pour $|u|=s$,
\begin{equation}
  \frac1n \binom{n-1}{s}^{-1} \;=\; \frac1n \cdot \frac{s!\,(n-1-s)!}{(n-1)!}
  \;=\; \frac{s!\,(n-s-1)!}{n!},
  \label{eq:poids-shapley}
\end{equation}
de sorte que $\varphi_i$ est exactement la moyenne, sur les $n!$ \textbf{permutations} de
$\{1,\dots,n\}$, de la contribution marginale de $i$ à l'ensemble de ses prédécesseurs. En effet,
un sous-ensemble $u$ de cardinal $s$ ne contenant pas $i$ est l'ensemble des prédécesseurs de $i$
dans exactement $s!\,(n-s-1)!$ permutations : $s!$ ordres pour les éléments de $u$, puis
$(n-s-1)!$ pour ceux qui suivent $i$. En notant $\Pi$ l'ensemble des permutations et $P_i(\pi)$
les prédécesseurs de $i$ dans $\pi$,
\[
  \varphi_i \;=\; \frac{1}{n!} \sum_{\pi \in \Pi}
              \bigl[v\bigl(P_i(\pi) \cup \{i\}\bigr) - v\bigl(P_i(\pi)\bigr)\bigr].
\]
Sommons alors sur $i$ \emph{à permutation fixée}. Si $\pi = (\pi(1),\dots,\pi(n))$, les
ensembles de prédécesseurs emboîtés donnent
$P_{\pi(j+1)}(\pi) = P_{\pi(j)}(\pi) \cup \{\pi(j)\}$, si bien que la somme des contributions
marginales le long de $\pi$ est \textbf{télescopique} :
\[
  \sum_{i=1}^{n} \bigl[v\bigl(P_i(\pi)\cup\{i\}\bigr)-v\bigl(P_i(\pi)\bigr)\bigr]
  \;=\; \sum_{j=1}^{n} \bigl[v\bigl(P_{\pi(j+1)}(\pi)\bigr)-v\bigl(P_{\pi(j)}(\pi)\bigr)\bigr]
  \;=\; v(\{1,\dots,n\}) - v(\emptyset),
\]
et cette valeur ne dépend pas de $\pi$. En moyennant sur les $n!$ permutations,
$\sum_i \varphi_i = v(\{1,\dots,n\})$, puisque $v(\emptyset)=0$.
\end{proof}

\paragraph{Ce que l'efficience garantit, et ce qu'elle ne garantit pas.}
Elle garantit que l'attribution est une partition \emph{exacte} du surcoût, quelle que soit
l'interaction et quel que soit son signe : rien n'est perdu ni dupliqué. Elle ne garantit
\textbf{pas} que la part d'un joueur soit une prédiction de ce qu'on obtiendrait en l'écartant
seul, laquelle est la contribution \emph{marginale} et ne somme pas au total. C'est la distinction
que le chapitre~\ref{chap:resultats} tient entre les trois lectures d'un même canal, et l'annexe
des adaptations par pilier en donne l'identité correspondante pour chaque colonne. En particulier,
une part de Shapley est une \emph{convention} d'attribution défendable, non un budget de
remédiation : deux des quatre canaux du modèle sont bornés et non calibrés.

\paragraph{Pourquoi aucun échantillonnage n'est nécessaire ici.}
La réécriture~\eqref{eq:poids-shapley} est aussi ce qui rend le calcul praticable : elle ramène
l'évaluation à une moyenne sur les permutations, forme sur laquelle reposent les algorithmes
d'approximation en temps polynomial \citep{CastroGomezTejada2009, SongNelsonStaum2016},
indispensables dès que la dimension croît. Avec cinq piliers, les $5!$ permutations et les $2^5$
valeurs de $v$ sont énumérées exhaustivement : les valeurs publiées sont donc exactes au bruit de
simulation de $v$ près, et aucune erreur d'approximation combinatoire ne s'y ajoute.

